% FILE: 02_lineage_and_context.tex
% Project: research-open-ontologies
% Role: open-ontology-research-and-rdf-authority

\section{Lineage and Context}
\label{sec:research-open-ontologies-lineage}

\subsection{2016 Language-Model Lesson}
\label{sec:research-open-ontologies-lm-lesson}

The 2016 language-model experiment established that local text imitation does not conserve
consequence. A system that predicts the next token in a process description does not thereby
understand the process. Predictions are statistically plausible but causally ungrounded: a
predicted process step has no claim on the event log, no receipt tying it to an observed
execution, and no typed witness connecting it to a formal theory of process.

The open-ontologies project is the direct architectural response to this lesson at the
knowledge representation layer. Rather than relying on statistical patterns over text, it
requires that every process intelligence concept --- every class, property, individual, and
relationship --- be expressed as a typed RDF triple in a public standard vocabulary or in a
local extension namespace that is explicitly declared and referenced. A concept that cannot
be expressed as a triple cannot be queried, cannot be validated by SHACL, and cannot serve
as the input to a receipted manufacturing rule. This is not a style preference; it is an
architectural enforcement of the consequence-conservation requirement exposed by the 2016
experiment.

\subsection{Enterprise Process Gap}
\label{sec:research-open-ontologies-enterprise-gap}

The enterprise process gap is the observation that organizations generate vast quantities of
event data but lack the vocabulary infrastructure to make auditable claims about what that
data proves. Fitness scores are reported in dashboards without naming the conformance algorithm,
the reference model, or the event log used. M\&A due diligence documents claim ``process
efficiency'' without specifying what process model was verified, at what fitness threshold,
against which event log, with what receipt.

The open-ontologies project addresses this gap by defining a shared vocabulary that makes
such claims expressible as structured assertions. \texttt{wasm4pm:ConformanceVerdict}
requires \texttt{fitness >= 0.95} and \texttt{precision >= 0.90} backed by a
\texttt{BLAKE3 CryptographicReceipt}. \texttt{ma:BoardClaim} is an OWL class hierarchy
with subclasses \texttt{SynergyProjection}, \texttt{OperationalDebtClaim},
\texttt{IntegrationRiskClaim}, \texttt{ProcessAssetClaim}, and \texttt{ControlClaim} ---
each requiring a named receipt type and a named witness marker before the claim is
board-admissible. Without this ontology layer, the claims remain informal and therefore
inadmissible under the program COVENANT.

\subsection{Relationship to the Chatman Equation}
\label{sec:research-open-ontologies-chatman-equation}

The Chatman Equation, $A = \mu(O^*)$, defines process intelligence as the composition of
admissible evidence ($A$) with the manufacturing function ($\mu$) applied to the open-source
manufacturing platform ($O^*$). The open-ontologies project provides the type system that
makes ``admissible'' a formal predicate rather than a subjective judgment.

Specifically: evidence $A$ is admissible in the sense of the Chatman Equation if and only if
it can be expressed as a collection of RDF triples conforming to the shapes defined by the
SHACL constraints in the corpus, grounded in a named \texttt{compat:WitnessMarker} from the
witness lattice, and accompanied by a BLAKE3 receipt. The open-ontologies corpus therefore
instantiates the admissibility predicate of the equation. It is the formal interface between
raw execution evidence (OTel traces, event logs) and the manufacturing function $\mu$ that
transforms evidence into board-admissible artifacts.

The federated dependency graph's \texttt{gate\_pass: true} status (Phase~4, 2026-06-01) is
therefore not merely a technical validation result; it is the certification that the
admissibility predicate of the Chatman Equation is currently satisfiable for the ontology corpus.

\subsection{Relationship to Receipts and Replay}
\label{sec:research-open-ontologies-receipts-replay}

The receipt chain is the cryptographic enforcement mechanism for the admissibility predicate.
Seven canonical receipts in \path{receipts/RECEIPT\_REGISTRY.md} certify the major research
program milestones: PAPER\_CANON, PM4PY\_ORACLE, WASM4PM\_GAP, LIFECYCLE, MA, STANDARDS, and
ADVERSARIAL. Each receipt names the operator that produced it, the named witness it was
produced under, and the open items that remain.

The open-ontologies project sits at the junction between the receipt chain and the ontology
corpus. The ggen manufacturing pipeline (declared in \path{ggen/ggen.toml}) uses BLAKE3
checksums as its \texttt{checksum\_algorithm} and \texttt{cryptographic-chain} as its
\texttt{proof\_format}. This means every artifact manufactured by a SPARQL-to-Tera rule is
itself a node in a receipted chain that can be replayed: given the same ontology graph and
the same SPARQL query, the same artifact must be reproducible with the same BLAKE3 hash.

Replay is the Van der Aalst Constitution's requirement applied to manufacturing: do not trust
the code path that claims to have manufactured an artifact; trust only the evidence that the
artifact can be re-derived from a conforming knowledge source. The open-ontologies project
provides that conforming knowledge source.
